\section{Motivation, Leadership and Entrepreneurship}
{\color{sub}\footnotesize 10 hrs \gap$\cdot$\gap asked in \mk{all 29 papers}
\gap$\cdot$\gap 3 units, 21 syllabus topics \gap$\cdot$\gap \mk{579 marks, 24.6\% of the archive}
$\rightarrow$ about 13 of every 80 marks: usually two full 8-mark questions.}

\vspace{3pt}\noindent
\colorbox{gdL}{\makebox[\dimexpr\textwidth-2\fboxsep][l]{%
  \textcolor{gd}{\bfseries\footnotesize READ THIS FIRST}}}
\par\nobreak\vspace{3pt}

Asked in every paper, like Ch.\,1 and Ch.\,2. Three blocks, each with a fixed pattern:
\begin{enumerate}[label=\arabic*., leftmargin=1.4em, itemsep=1pt]
  \item \textbf{Motivation} = ``define motivation'' (2--3 marks) $+$ one named theory
    (5--6 marks). Five theories cover it: \textbf{Maslow} (9 papers), \textbf{Herzberg} (10),
    \textbf{Vroom} (7), \textbf{McGregor X-Y} (6), \textbf{ERG} as a comparison (6). Learn one
    figure or table per theory. \textbf{Content} theories (Maslow, ERG, McClelland,
    Herzberg) say \textit{what} motivates; \textbf{process} theory (Vroom) says \textit{how}.
  \item \textbf{Leadership} = qualities of a good leader (12) or styles (11) $+$ an opinion:
    ``which style suits an engineering project / Nepal / a startup?''. Answer: situational,
    leaning democratic, with a reason for each context ($\S$3.3).
  \item \textbf{Entrepreneurship} = definition + characteristics (8) or ``why promote it in
    Nepal'' (7). Always tie the second half to Nepal: jobs, remittance dependence, IT.
\end{enumerate}
Tier chips count \textbf{papers out of 29}: \tS{n} 8+ \gap \tF{n} 4--7 \gap \tP{n} 1--3
inside an 8-mark question \gap \tL{n} 1--3, short note only \gap \tN{} never set.
Bold year = Regular exam.

\hr

%=======================================================================
\T{3.1 Motivation: Needs, Types, Techniques}

\Q{Define motivation. What are human needs, and how is a need used to motivate? Types of motivation}{\m{2}\m{3}\m{4}}
{\color{sub}\footnotesize \tS{16} \yr{\textbf{82 Bh}, \textbf{79 Bh}, 79 Ba, \textbf{76 Ch}, 76 Ash, \textbf{75 Ch}, \textbf{74 Ch}, 74 Ash, \textbf{73 Ch}, 72 Ka, \textbf{71 Ch}, 71 Shr, \textbf{70 Ch}, 70 Asa, \textbf{69 Ch}, \textbf{68 Ch}} \gap$\cdot$\gap the 2--3 mark opener to nearly every theory question}

\sbs{%
\begin{itemize}
  \item From Latin \textit{movere}, ``to move''; \textbf{motive} = need, want, drive within a
    person.
  \item \textbf{Definition}: the process of \textbf{stimulating people to act} toward goals
    by creating conditions in which \textbf{their needs and the organization's goals are
    satisfied together}.
  \item Forces within a person that decide the \textbf{direction, level and persistence} of
    effort at work.
  \item Two aspects: knowing human needs; inspiring people to do their best for the
    enterprise.
\end{itemize}
\lead{Human needs}
\begin{itemize}
  \item \textbf{Primary} (physiological, safety): food, shelter, security.
  \item \textbf{Secondary} (social, esteem, self-actualization): belonging, status, growth.
  \item Needs are unlimited, change over time, differ between people, and several operate at
    once.
\end{itemize}
}{%
\lead{Types of motivation}
\begin{itemize}
  \item \textbf{Intrinsic}: from within; interest, enjoyment, sense of achievement in the task
    itself. ``I want to complete this task.'' Lasting.
  \item \textbf{Extrinsic}: from outside; pay, bonus, promotion, fear of penalty.
    ``Complete it and be rewarded.''
  \item \textbf{Positive}: reward, praise, incentive (carrot).
  \item \textbf{Negative}: fear, punishment, threat of dismissal (stick).
  \item \textbf{Financial vs non-financial}.
  \item By drive (Adhikari): incentive, fear, achievement, growth, power, social.
\end{itemize}
}

\sbs{%
\lead{How a need motivates: the need-drive-goal cycle}
\begin{enumerate}
  \item \textbf{Unsatisfied need} $\rightarrow$ tension.
  \item Tension $\rightarrow$ \textbf{drive} (motive) to reduce it.
  \item Drive $\rightarrow$ \textbf{goal-directed behaviour} (work harder).
  \item Goal reached $\rightarrow$ \textbf{need satisfied}, tension falls.
  \item New need appears; cycle repeats.
\end{enumerate}
Manager's job: \textbf{link the employee's unmet need to the organization's goal}, so working
for the firm satisfies it. A satisfied need no longer motivates.}{%
\lead{Why motivation is complex}
\begin{itemize}
  \item Same need $\rightarrow$ different behaviour in different people.
  \item Needs change in kind and intensity.
  \item People cannot always express their needs.
  \item Many needs act at once.
\end{itemize}
\lead{Importance}
\begin{itemize}
  \item Higher productivity and quality; better use of resources.
  \item Low absenteeism and turnover; fewer grievances.
  \item Willing acceptance of change; good relations.
  \item Helps employees reach personal goals.
\end{itemize}}

\Q{Techniques of motivation; attitude, group and executive motivation; role of management}{\m{2}\m{3}\m{4}\m{5}\m{8}}
{\color{sub}\footnotesize \tS{9} \yr{82 Ba, \textbf{81 Bh}, 81 Ba, \textbf{76 Ch}, 75 Ash, 73 Shr, 72 Ka, 70 Asa, \textbf{69 Ch}}}

\sbs{%
\lead{Techniques}
\begin{enumerate}
  \item \textbf{Financial}: fair wages, bonus, incentives, profit sharing, ESOP.
  \item \textbf{Meaningful, challenging work}: job enrichment, rotation.
  \item \textbf{Clear targets} and measured performance (MBO).
  \item \textbf{Regular, specific feedback}.
  \item \textbf{Recognition and praise}: public appreciation, awards.
  \item \textbf{Participation} in decisions; \textbf{delegation}.
  \item \textbf{Training, career growth}, promotion.
  \item \textbf{Job security}, good working conditions, quality of work life.
  \item \textbf{Status}, healthy competition, flexible hours.
\end{enumerate}
\lead{Role of management in motivation}
\begin{itemize}
  \item Understand each employee's needs; people differ.
  \item Create an environment where personal and firm goals meet.
  \item Fair pay system; fulfil promised rewards.
  \item Give information, training, feedback; act on employee suggestions.
  \item Lead by example; make work interesting.
\end{itemize}}{%
\lead{Attitude (self) motivation}
\begin{itemize}
  \item How a person thinks and feels: self-confidence, belief, positive outlook.
  \item Falls with monotonous work, pushy boss, poor health, bad relations, money worries.
  \item Cure: interesting job, constructive thinking, build on strengths, adapt.
\end{itemize}
\lead{Group motivation}
\begin{itemize}
  \item Motivating a \textbf{team} as a unit; team productivity decides success.
  \item Needs fairness, honesty, loyalty, tolerance, shared responsibility, constructive
    criticism.
  \item Tools: team goals, group incentives, team building, social recognition.
\end{itemize}
\lead{Executive motivation}
\begin{itemize}
  \item Keeping \textbf{managerial talent} motivated and in the firm.
  \item Executives seek achievement, esteem, self-actualization more than money.
  \item Tools: challenging responsibility, autonomy, profit sharing, stock options, status
    perks, public recognition.
  \item Low when: poor boss, blocked advancement, poor pay, insecurity.
\end{itemize}
{\small \textbf{Difference} (75 Ash): attitude = individual's inner state; group = team as a
whole; executive = top talent, higher-order needs.}}

\penalty0\vspace{2pt}

%=======================================================================
\T{3.2 Motivation Theories}

\Q{Maslow's hierarchy of needs: features, satisfaction-progression, examples}{\m{3}\m{4}\m{5}}
{\color{sub}\footnotesize \tS{9} \yr{\textbf{82 Bh}, 82 Ba, \textbf{81 Bh}, \textbf{78 Bh}, 76 Ash, \textbf{75 Ch}, 74 Ash, \textbf{70 Ch}, \textbf{69 Ch}}}

\sbsr{0.50}{%
\begin{itemize}
  \item \textbf{Abraham Maslow}, 1943. Three propositions:
  \begin{itemize}
    \item Man is a \textbf{wanting being}: needs are endless.
    \item A \textbf{satisfied need does not motivate}; only unmet needs do.
    \item Needs form a \textbf{hierarchy}: lower must be largely met before higher ones
      motivate.
  \end{itemize}
  \item \textbf{Five levels} (bottom up), with what the firm gives:
  \begin{enumerate}
    \item \textbf{Physiological}: food, water, shelter, sleep $\rightarrow$ adequate wage,
      rest breaks.
    \item \textbf{Safety}: physical and economic security $\rightarrow$ safe conditions,
      job security, insurance, pension.
    \item \textbf{Social}: belonging, friendship $\rightarrow$ teams, informal groups,
      events.
    \item \textbf{Esteem}: self-respect, status, recognition $\rightarrow$ titles, awards,
      promotion.
    \item \textbf{Self-actualization}: realizing full potential $\rightarrow$ challenging,
      creative work, autonomy.
  \end{enumerate}
\end{itemize}}{%
\figT{c3_maslow.png}
{\color{sub}\scriptsize Maslow's pyramid. Lower needs must be met before higher ones
motivate. Adhikari notes, p197.}}

\lead{Satisfaction-progression process (75 Ch)}
\begin{itemize}
  \item When a need is \textbf{satisfied}, the person \textbf{progresses} to the next higher
    need, which now motivates.
  \item Eg a fresh graduate first wants a salary (physiological), then a permanent post
    (safety), then a good team (social), then a senior title (esteem), then to start
    something of their own (self-actualization).
\end{itemize}

\sbs{%
\lead{Features and limitations}
\begin{itemize}
  \item $+$ Simple, intuitive; widely used by managers.
  \item $-$ Order is not fixed for everyone (an artist may skip safety).
  \item $-$ Several needs act together.
  \item $-$ Hard to measure satisfaction; culture-bound (US study).
  \item $-$ Only one direction (no regression) $\rightarrow$ ERG fixes this.
\end{itemize}
\lead{Can Maslow explain Laxmi Prasad Devkota's tireless quest? (69 Ch)}
\begin{itemize}
  \item \textbf{Only partly}. Devkota wrote \textit{Muna Madan} and his epics while in
    poverty, debt and grief, and kept writing on his deathbed with cancer.
  \item Pure Maslow says unmet physiological and safety needs should have blocked
    self-actualization.
  \item So he \textbf{contradicts the strict order}: self-actualization (creative
    expression, intrinsic motivation) drove him even with lower needs unmet.
  \item Better explained by \textbf{ERG} (needs act together) and intrinsic motivation.
\end{itemize}}{%
\lead{Why some fall from grace while climbing the pyramid (78 Bh)}
\begin{itemize}
  \item Fraud, crime, suicide among the successful are the \textbf{irony} of the pyramid.
  \item \textbf{Esteem sought by any means}: status, wealth, power without ethics
    $\rightarrow$ corruption, fraud.
  \item \textbf{Frustration-regression} (ERG): blocked at the top, people fall back to
    greed for material things or to despair.
  \item \textbf{Personal power need} (McClelland) outgrows social good.
  \item Pressure, isolation, loss of social support $\rightarrow$ depression, suicide.
  \item Maslow's own top level includes \textbf{morality}: those who fall never truly reached
    self-actualization, they stalled at esteem.
  \item Remedy: values and ethics education, social support, fair systems, counselling.
\end{itemize}}

\Q{Maslow vs Alderfer's ERG; ERG and Herzberg compared with Maslow; satisfaction-progression vs frustration-regression}{\m{3}\m{4}\m{5}\m{6}}
{\color{sub}\footnotesize \tF{6} \yr{82 Ba, \textbf{78 Bh}, \textbf{75 Ch}, 71 Shr, \textbf{69 Ch}, \textbf{68 Ch}} \gap$\cdot$\gap satisfaction-progression vs frustration-regression as a short note in 75 Ch, 69 Ch}

\sbs{%
\lead{Alderfer's ERG theory (1969)}
\begin{itemize}
  \item Collapses Maslow's five into three:
  \begin{itemize}
    \item \textbf{Existence}: physiological + safety.
    \item \textbf{Relatedness}: social + external esteem.
    \item \textbf{Growth}: internal esteem + self-actualization.
  \end{itemize}
  \item \textbf{Satisfaction-progression}: satisfied lower need $\rightarrow$ move up (as
    Maslow).
  \item \textbf{Frustration-regression}: if a higher need is \textbf{blocked}, the person
    \textbf{falls back} and wants more of the lower need. Eg an engineer denied promotion
    (growth) demands higher pay (existence).
  \item Several needs can motivate \textbf{at the same time}.
\end{itemize}}{%
{\small
\begin{tabularx}{\linewidth}{@{}L{1.5cm}XX@{}}
\toprule
\hrow \thead{Basis} & \thead{Maslow} & \thead{Alderfer ERG} \\
\midrule
Levels & 5 & 3 \\
Order & Strict, step by step & Flexible \\
Needs at once & One dominant & Several together \\
Direction & Up only (progression) & Up and down (regression) \\
Basis & Clinical observation & Empirical research \\
\bottomrule
\end{tabularx}}
\vspace{3pt}
\lead{Herzberg against Maslow and ERG}
\begin{itemize}
  \item Herzberg's \textbf{hygiene factors} = Maslow's lower needs = ERG's existence and
    relatedness.
  \item Herzberg's \textbf{motivators} = Maslow's esteem + self-actualization = ERG's growth.
  \item Difference: Maslow and ERG say every need motivates when unmet; Herzberg says lower
    needs only \textbf{prevent dissatisfaction}, never motivate.
  \item Maslow and ERG: general human needs. Herzberg: \textbf{job} factors.
\end{itemize}}

\Q{McGregor's Theory X and Theory Y}{\m{5}\m{6}\m{8}}
{\color{sub}\footnotesize \tF{6} \yr{\textbf{81 Bh}, 81 Ba, 79 Ba, 73 Shr, \textbf{72 Ch}, \textbf{65 Shr}}}
\sbs{%
\begin{itemize}
  \item \textbf{Douglas McGregor}, \textit{The Human Side of Enterprise} (1960): two sets of
    \textbf{assumptions managers hold about workers}, which shape how they manage.
\end{itemize}
\lead{Theory X (traditional)}
\begin{itemize}
  \item Average person dislikes work, avoids it.
  \item No ambition, avoids responsibility, prefers to be led.
  \item Self-centred, indifferent to organizational goals.
  \item Resists change; motivated only by money and fear.
  \item $\rightarrow$ Manager must \textbf{coerce, control, direct, threaten}: autocratic
    style, close supervision.
\end{itemize}}{%
\lead{Theory Y (modern)}
\begin{itemize}
  \item Work is as natural as play or rest.
  \item People exercise \textbf{self-direction and self-control} when committed.
  \item Commitment comes from rewards of achievement.
  \item They \textbf{seek responsibility}; creativity is widespread.
  \item $\rightarrow$ Manager \textbf{delegates, involves, enriches jobs}: democratic style.
\end{itemize}
\lead{X vs Y against the other theories (72 Ch)}
\begin{itemize}
  \item X assumes people live at Maslow's \textbf{lower needs} / Herzberg's hygiene level;
    fits fear-and-punishment theory.
  \item Y assumes \textbf{higher needs} and Herzberg's motivators drive people.
  \item McGregor favoured Y; modern view: situation decides (contingency).
\end{itemize}}

\Q{Herzberg's motivation-hygiene (two-factor) theory: where motivation at work comes from}{\m{4}\m{5}\m{6}}
{\color{sub}\footnotesize \tS{10} \yr{82 Ba, \textbf{80 Bh}, \textbf{79 Bh}, \textbf{78 Bh}, \textbf{74 Ch}, \textbf{73 Ch}, \textbf{71 Ch}, \textbf{68 Ba}, \textbf{67 Asa}, \textbf{66 Bh}}}
\sbs{%
\begin{itemize}
  \item \textbf{Frederick Herzberg}, 1959: interviewed about 200 engineers and accountants in
    Pittsburgh on times they felt very good or very bad at work.
  \item Finding: \textbf{satisfaction and dissatisfaction come from different factors}. The
    opposite of satisfaction is \textbf{no satisfaction}, not dissatisfaction.
  \item Man has two sets of needs: as an animal, to \textbf{avoid pain}; as a human, to
    \textbf{grow psychologically}.
\end{itemize}
\lead{Hygiene (maintenance) factors: job context}
\begin{itemize}
  \item Company policy and administration; supervision.
  \item Salary; job security; working conditions.
  \item Relations with superiors, peers, subordinates; status.
  \item Absent $\rightarrow$ \textbf{dissatisfaction}. Present $\rightarrow$ only \textbf{no
    dissatisfaction}, not motivation.
\end{itemize}}{%
\lead{Motivators: job content}
\begin{itemize}
  \item Achievement; recognition.
  \item The work itself (interesting, challenging).
  \item Responsibility; advancement; growth.
  \item Present $\rightarrow$ \textbf{satisfaction and motivation}. Absent $\rightarrow$ no
    satisfaction (but not dissatisfaction).
\end{itemize}
\lead{So where does motivation at work come from? (80 Bh)}
\begin{itemize}
  \item From the \textbf{job itself}: achievement, recognition, responsibility, growth.
    Pay and conditions only remove complaints.
  \item Managers: first fix hygiene (else unhappiness), then \textbf{enrich jobs} to motivate.
\end{itemize}
\lead{Criticism}
\begin{itemize}
  \item Small, professional sample; people credit themselves for good events and blame
    others for bad (method bias).
  \item Pay can be a motivator for low-income workers.
\end{itemize}}

\Q{Vroom's expectancy (VIE / valency) theory; why it is a process theory; importance of reward}{\m{4}\m{5}\m{8}}
{\color{sub}\footnotesize \tF{7} \yr{\textbf{82 Bh}, \textbf{80 Bh}, 80 Ba, \textbf{78 Bh}, \textbf{76 Ch}, \textbf{74 Ch}, 70 Asa} \gap$\cdot$\gap short note in 82 Bh}
\sbsr{0.50}{%
\begin{itemize}
  \item \textbf{Victor Vroom}, 1964. A person chooses the behaviour they \textbf{expect} to
    bring the outcome they \textbf{value}.
  \item \[ \text{Motivation } M = E \times I \times V \]
  \item \textbf{Expectancy (E)}: belief that \textbf{effort $\rightarrow$ performance}
    (0 to 1). ``If I try, can I do it?''
  \item \textbf{Instrumentality (I)}: belief that \textbf{performance $\rightarrow$
    reward}. ``If I do it, will I be rewarded?''
  \item \textbf{Valence (V)}: \textbf{value} the person places on that reward. ``Do I want
    it?''
  \item Product: if \textbf{any one is zero}, motivation is zero.
  \item Eg an engineer offered a Dubai posting (V high) for finishing a project early: no
    effort if the deadline is impossible (E low) or if the boss rarely keeps promises (I low).
\end{itemize}}{%
\figT{c3_vroom.png}
{\color{sub}\scriptsize Effort $\rightarrow$ performance $\rightarrow$ outcomes, and the belief
at each link. Adhikari notes, p219.}
\lead{Why it is a process theory (80 Bh)}
\begin{itemize}
  \item It does not list \textit{which} needs motivate (content theories do).
  \item It explains \textbf{how} a person \textbf{decides}: a mental calculation of E, I, V
    among alternatives.
\end{itemize}
\lead{Importance of reward (74 Ch)}
\begin{itemize}
  \item Reward sits in two of three terms: it must be \textbf{valued} (V) and \textbf{reliably
    linked} to performance (I).
  \item Managers: offer rewards people want; keep promises; make targets reachable by
    training and resources (E).
\end{itemize}}

\creamq{McClelland's theory of learned needs}{}
{\color{sub}\footnotesize \tN}
\begin{itemize}
  \item David McClelland (1961): needs are \textbf{learned} from life experience.
  \item \textbf{Achievement (nAch)}: excel, moderate risks, want feedback, personal
    responsibility $\rightarrow$ entrepreneurs.
  \item \textbf{Power (nPow)}: influence others; personal or social (institutional) power
    $\rightarrow$ managers.
  \item \textbf{Affiliation (nAff)}: be liked, avoid conflict $\rightarrow$ team roles.
\end{itemize}

\creamq{Fear and punishment theory}{\m{4}}
{\color{sub}\footnotesize \tL{1} \yr{\textbf{82 Bh}} \gap$\cdot$\gap short note}
\begin{itemize}
  \item People work to \textbf{avoid penalties}: warnings, fines, demotion, dismissal. Linked to
    Theory X and Skinner's reinforcement (punishment, negative reinforcement).
  \item \textbf{Deterrence}: specific (corrects the offender), general (warns others).
  \item Works short term for discipline, safety rules.
  \item Causes resentment, fear, minimum effort, turnover; use only as a last resort.
\end{itemize}

\penalty0\vspace{2pt}

%=======================================================================
\T{3.3 Leadership}

\Q{Define leadership. Why is it needed? Qualities of a good leader}{\m{3}\m{4}\m{5}\m{8}}
{\color{sub}\footnotesize \tS{12} \yr{\textbf{81 Bh}, 81 Ba, \textbf{80 Bh}, \textbf{78 Bh}, \textbf{76 Ch}, 75 Ash, 74 Ash, \textbf{73 Ch}, 72 Ka, \textbf{70 Ch}, 70 Asa, \textbf{68 Ch}}}
\sbs{%
\begin{itemize}
  \item \textbf{Leadership}: the process of \textbf{influencing people} to work willingly
    toward group or organizational goals, and providing an environment for it.
  \item Koontz: ``the art or process of influencing people so that they strive willingly and
    enthusiastically toward the achievement of group goals.''
\end{itemize}
\lead{Why needed}
\begin{itemize}
  \item Gives \textbf{vision and direction}.
  \item \textbf{Motivates}, builds morale and confidence.
  \item Coordinates effort; resolves conflict.
  \item Leads \textbf{change}; builds culture and teamwork.
  \item Represents the group; develops future leaders.
\end{itemize}
}{%
\lead{Qualities of a good leader}
\begin{enumerate}
  \item \textbf{Vision}: knows where to go, sets achievable goals.
  \item \textbf{Honesty, integrity}: same inside and out; trust.
  \item \textbf{Self-confidence}, courage, decisiveness.
  \item \textbf{Communication}: clear speaking, writing, \textbf{listening}.
  \item \textbf{Human skill}: respect, fairness, empathy; credit to others.
  \item \textbf{Initiative} and \textbf{innovation}.
  \item \textbf{Knowledge}, intelligence, judgement.
  \item \textbf{Discipline}, commitment to excellence; leads by example.
  \item \textbf{Emotional stability}, patience, sense of humour.
  \item \textbf{Calculated risk-taking}; learns from mistakes.
  \item \textbf{Charisma}, passion.
  \item \textbf{Adaptability}: changes style with the situation.
\end{enumerate}}

\lead{What sets a leader apart from regular employees (80 Bh)}
\begin{itemize}
  \item Vision beyond own job; takes initiative unasked.
  \item Owns results, including failures.
  \item Influences without formal authority; others follow willingly.
  \item Develops others; thinks of the team first.
\end{itemize}

\Q{Leadership styles (autocratic, democratic, free-rein); which style suits an engineering project, industry, Nepal, a startup?}{\m{2}\m{3}\m{4}\m{5}\m{6}\m{7}\m{8}}
{\color{sub}\footnotesize \tS{11} \yr{82 Ba, 80 Ba, \textbf{78 Bh}, \textbf{75 Ch}, 74 Ash, \textbf{73 Ch}, \textbf{72 Ch}, \textbf{68 Ch}, \textbf{68 Ba}, \textbf{67 Asa}, \textbf{65 Shr}} \gap$\cdot$\gap ``leadership style'' short note in 68 Ba \gap$\cdot$\gap \mk{opinion half every time}}

\begin{itemize}
  \item \textbf{Style} = the behaviour pattern a leader adopts; depends on the leader's
    personality, followers, task, tradition, time available.
\end{itemize}

\sbsr{0.55}{%
\begin{itemize}
  \item \textbf{Autocratic (authoritarian)}:
  \begin{itemize}
    \item Leader decides alone, gives orders, expects obedience; one-way communication.
    \item $+$ quick decisions, clear control; suits emergencies, unskilled or new staff,
      military.
    \item $-$ low morale, no creativity, dependence on one person, resentment.
  \end{itemize}
  \item \textbf{Democratic (participative)}:
  \begin{itemize}
    \item Leader consults, decides with the group; praise and blame shared; two-way.
    \item $+$ better decisions, commitment, morale, creativity; usually most productive.
    \item $-$ slow; needs capable, willing followers.
  \end{itemize}
  \item \textbf{Laissez-faire (free-rein)}:
  \begin{itemize}
    \item Leader gives goals and freedom; group decides and self-controls.
    \item $+$ suits experts, researchers; high autonomy.
    \item $-$ no direction, chaos with inexperienced staff.
  \end{itemize}
  \item Others: \textbf{paternalistic} (fatherly care), \textbf{bureaucratic} (by rules),
    \textbf{transformational} (inspires change), \textbf{transactional} (reward-punishment).
\end{itemize}}{%
\figT{c3_styles.png}
{\color{sub}\scriptsize Power, decision making and communication under the three styles.
Adhikari notes, p238.}
\lead{Participative management (65 Shr)}
\begin{itemize}
  \item Employees share in decisions: suggestion schemes, quality circles, work committees,
    worker directors.
  \item Importance: commitment, better ideas, fewer disputes, job satisfaction.
\end{itemize}}

\sbs{%
\lead{Which style for a (modern) engineering project? (78 Bh, 74 Ash)}
\begin{itemize}
  \item \textbf{Democratic / participative} as the base: engineers are skilled
    professionals; design choices gain from many views; commitment matters over months.
  \item \textbf{Autocratic} moments: safety, deadlines, site emergencies.
  \item \textbf{Free-rein} for expert sub-teams (design, R\&D).
  \item $\rightarrow$ \textbf{Situational leadership}, democratic by default.
\end{itemize}
\lead{Most effective style in an industrial organization (80 Ba, 68 Ba, 67 Asa)}
\begin{itemize}
  \item Determine it by: task structure, followers' skill and maturity, leader's power,
    time pressure (Fiedler, Hersey-Blanchard).
  \item Shop floor, routine work, new workers $\rightarrow$ directive; skilled staff
    $\rightarrow$ participative; experts $\rightarrow$ delegating.
  \item Most effective overall: \textbf{team (9,9) style} on Blake-Mouton's grid.
\end{itemize}}{%
\lead{Most efficient in Nepal? (73 Ch)}
\begin{itemize}
  \item Tradition is autocratic and paternalistic (hierarchy, respect for seniority).
  \item Educated young workforce, migration of talent $\rightarrow$ \textbf{democratic with
    a paternalistic touch} works best: consult, care for welfare, decide firmly.
  \item Government offices need less bureaucracy and more participation.
\end{itemize}
\lead{Nepalese start-ups (82 Ba)}
\begin{itemize}
  \item Small teams of young, skilled people, fast change, scarce funds.
  \item Best: \textbf{democratic + transformational}: shared vision, open ideas, quick
    consensus, ownership (ESOP), freedom for experts.
  \item Founder must switch to decisive (autocratic) mode in a cash or deadline crisis.
\end{itemize}}

\Q{Blake and Mouton's managerial grid: styles, examples, significance for growth}{\m{4}\m{5}\m{8}}
{\color{sub}\footnotesize \tF{6} \yr{\textbf{82 Bh}, \textbf{80 Bh}, 79 Ba, 76 Ash, 75 Ash, \textbf{70 Ch}}}
\sbsr{0.48}{%
\begin{itemize}
  \item Robert Blake and Jane Mouton (1964). Two axes, each 1 (low) to 9 (high):
    \textbf{concern for production} (x) and \textbf{concern for people} (y).
  \item Five key styles:
  \begin{itemize}
    \item \textbf{(1,1) Impoverished}: minimum effort for either; avoids blame. Eg a manager
      waiting for retirement.
    \item \textbf{(1,9) Country club}: comfort and friendliness, weak output. Eg a boss who
      never enforces deadlines.
    \item \textbf{(9,1) Authority-compliance (task)}: output through rules and pressure,
      people ignored. Eg a sweatshop supervisor.
    \item \textbf{(5,5) Middle of the road}: balance, adequate results, neither fully met.
    \item \textbf{(9,9) Team}: high on both; commitment, trust, shared stake. Eg an agile
      software team lead.
  \end{itemize}
\end{itemize}}{%
\figT{c3_grid.png}
{\color{sub}\scriptsize The leadership grid: concern for production vs concern for people.
Adhikari notes, p241.}
\lead{Significance for organizational growth (76 Ash)}
\begin{itemize}
  \item Diagnoses a manager's present style objectively.
  \item Shows the target: move toward \textbf{(9,9)}.
  \item Basis for leadership training (grid seminars).
  \item Shows people and production are not a trade-off: both together raise productivity
    and morale $\rightarrow$ growth.
\end{itemize}}

\Q{Leadership approaches / theories: trait, behavioural, contingency (situational), integrated}{\m{4}\m{5}\m{8}}
{\color{sub}\footnotesize \tF{7} \yr{79 Ba, \textbf{76 Ch}, \textbf{75 Ch}, 75 Ash, \textbf{69 Ch}, \textbf{66 Bh}, \textbf{65 Shr}} \gap$\cdot$\gap contingency approach of leadership as a short note in 66 Bh}
\sbs{%
\lead{Trait approach}
\begin{itemize}
  \item Leaders are \textbf{born, not made} (Great Man theory).
  \item Identifies traits of effective leaders: drive, leadership motivation, honesty,
    self-confidence, intelligence, business knowledge, self-monitoring.
  \item $-$ No trait list fits all leaders; ignores followers and situation.
\end{itemize}
\lead{Behavioural approach}
\begin{itemize}
  \item Leaders are \textbf{made}: what they \textbf{do} matters, and can be learned.
  \item Two independent behaviours: \textbf{task-oriented} (assign, set standards, push
    deadlines) and \textbf{people-oriented} (trust, respect, welfare). Ohio State and
    Michigan studies; Blake-Mouton grid.
  \item Effective leaders score high on both.
\end{itemize}}{%
\lead{Contingency / situational approach}
\begin{itemize}
  \item \textbf{No one best style}: it depends on leader, followers and situation.
  \item \textbf{Fiedler}: task- or relationship-motivated leader fits situations of
    different favourableness (leader-member relations, task structure, position power).
  \item \textbf{House's path-goal}: leader chooses \textbf{directive, supportive,
    participative or achievement-oriented} behaviour to clear the path to goals.
  \item \textbf{Hersey-Blanchard}: telling, selling, participating, delegating as follower
    maturity rises.
\end{itemize}
\lead{Integrated approach}
\begin{itemize}
  \item Combines \textbf{leader} (traits, behaviour, experience), \textbf{followers} (ability,
    maturity) and \textbf{situation} (structure, technology, objectives, environment) to
    choose the effective style. Most realistic.
\end{itemize}}

\Q{Leader vs manager; boss vs leader; ``a reader is a leader''; challenges for a modern leader; you as a future leader}{\m{2}\m{3}\m{4}\m{5}\m{8}}
{\color{sub}\footnotesize \tF{4} \yr{\textbf{75 Ch}, 75 Ash, \textbf{73 Ch}, \textbf{71 Ch}}}
\sbs{%
{\small
\begin{tabularx}{\linewidth}{@{}L{1.5cm}XX@{}}
\toprule
\hrow \thead{Basis} & \thead{Manager} & \thead{Leader} \\
\midrule
Power from & Position (authority) & Personal influence \\
Focus & Things, systems, control & People, vision, inspiration \\
Motto & Do things right & Do the right things \\
Work & Plan, organize, budget & Set direction, align, motivate \\
Change & Maintains order & Drives change \\
Followers & Subordinates & Followers \\
\bottomrule
\end{tabularx}}
A good manager should be a leader; a leader need not be a manager.
\lead{Boss vs leader (73 Ch)}
\begin{itemize}
  \item \textbf{Boss} = autocratic style: commands, uses fear, takes credit, blames, says ``I''.
  \item \textbf{Leader} = democratic / transformational style: guides, inspires, shares
    credit, owns failure, says ``we''.
\end{itemize}}{%
\lead{``A reader is a leader'' (71 Ch)}
\begin{itemize}
  \item Reading builds \textbf{knowledge} (expert power), \textbf{vision}, \textbf{judgement}
    from others' experience, and \textbf{communication}.
  \item Informed leaders can switch styles wisely: consult as a democrat, decide as an
    autocrat when needed.
  \item Eg Gandhi, Mandela, Gates: lifelong readers.
\end{itemize}
\lead{Challenges for a modern leader (75 Ch)}
\begin{itemize}
  \item Rapid technological change; remote and diverse teams.
  \item Talent retention (migration in Nepal); ethics, transparency.
  \item Uncertainty (pandemic, politics); information overload.
\end{itemize}
\lead{Me as a future leader (75 Ash, 75 Ch)}
\begin{itemize}
  \item Start as a technically strong engineer; build human skills.
  \item Lead by example, listen, share credit, keep learning.
  \item Aim for a (9,9) team style, adapting to the situation.
\end{itemize}}

\creamq{Authority and power}{\m{5}}
{\color{sub}\footnotesize \tP{1} \yr{\textbf{66 Bh}}}
\begin{itemize}
  \item \textbf{Authority}: \textbf{legitimate right} to command, from the position; flows
    downward; accepted as rightful.
  \item \textbf{Power}: \textbf{ability to influence} behaviour, from any source; can exist
    without authority.
  \item French and Raven's sources of power: \textbf{legitimate} (position), \textbf{reward},
    \textbf{coercive} (punish), \textbf{expert} (knowledge), \textbf{referent} (personal
    charisma). Leaders rely most on expert and referent power.
\end{itemize}

\penalty0\vspace{2pt}

%=======================================================================
\T{3.4 Entrepreneurship}

\Q{Entrepreneurship: meaning, development, characteristics of an entrepreneur}{\m{3}\m{4}\m{5}\m{8}}
{\color{sub}\footnotesize \tS{8} \yr{\textbf{79 Bh}, \textbf{76 Ch}, 75 Ash, 74 Ash, \textbf{72 Ch}, 71 Shr, 70 Asa, \textbf{69 Ch}} \gap$\cdot$\gap entrepreneurial characteristics as a short note in 79 Bh}
\begin{itemize}
  \item From French \textit{entreprendre}, ``to undertake''.
  \item \textbf{Entrepreneurship}: the capacity and willingness to \textbf{spot an
    opportunity, organize resources and run a new venture, bearing its risk}, for profit and
    value.
  \item \textbf{Schumpeter}: entrepreneurs are \textbf{innovators}: new products, new
    methods, new markets, new sources of supply, new organization.
  \item Functions: see opportunity, organize enterprise, raise capital, hire labour, arrange
    materials, select managers, bear risk.
  \item \textbf{McClelland}: a nation's growth depends on its people's need for achievement
    more than on natural resources.
\end{itemize}

\sbs{%
\lead{Entrepreneurship development}
\begin{itemize}
  \item Entrepreneurs are \textbf{made}: Entrepreneurship Development Programmes (EDPs).
  \item Steps: identify potential people $\rightarrow$ build capability (motivation,
    training) $\rightarrow$ help choose a viable project $\rightarrow$ teach management,
    finance $\rightarrow$ link to finance and infrastructure.
  \item Factors influencing it: need for achievement, independence, family background,
    creativity, flexibility, environment.
\end{itemize}}{%
\lead{Characteristics of an entrepreneur}
\begin{enumerate}
  \item \textbf{Initiative} and self-starting.
  \item \textbf{Opportunity seeking}; foresight.
  \item \textbf{Calculated (moderate) risk-taking}: not a gambler.
  \item \textbf{Creativity and innovation}.
  \item \textbf{High need for achievement}; competitive.
  \item \textbf{Self-confidence}, optimism.
  \item \textbf{Commitment, persistence}: learns from failure.
  \item \textbf{Decision making}; administrative ability.
  \item \textbf{People skills}, communication, leadership.
  \item \textbf{Technical and business knowledge}.
  \item \textbf{Hard work}, discipline, strong work ethic.
  \item \textbf{Flexibility}; handles uncertainty.
\end{enumerate}}

\Q{Steps for establishing a small scale unit}{\m{4}\m{5}}
{\color{sub}\footnotesize \tP{3} \yr{\textbf{81 Bh}, \textbf{79 Bh}, 71 Shr}}
\sbs{%
\begin{enumerate}
  \item \textbf{Self-assessment}: will, skills, risk appetite.
  \item \textbf{Product / service identification}: market research, competition, demand.
  \item \textbf{Feasibility study} and \textbf{project report}: technical, market, financial.
  \item \textbf{Form of ownership}: sole, partnership, private limited.
  \item \textbf{Location}: raw materials, market, transport, labour, power, water, waste
    disposal.
  \item \textbf{Registration and approvals}: Department / Office of Cottage and Small
    Industries (Nepal), PAN / VAT, local ward, environmental clearance.
\end{enumerate}}{%
\begin{enumerate}[start=7]
  \item \textbf{Finance}: fixed capital (land, building, machines) + working capital; own
    funds, bank loan, government subsidized loans.
  \item \textbf{Land, building, machinery}; utility connections.
  \item \textbf{Recruit and train} manpower.
  \item \textbf{Raw material} arrangement; \textbf{production} layout, trial run, quality
    control.
  \item \textbf{Marketing}: pricing, distribution, promotion.
  \item \textbf{Monitor and improve} the plan at every stage.
\end{enumerate}}

\Q{Need for promotion of entrepreneurship; its role in economic development; importance for Nepal; conducive environment and law enforcement}{\m{4}\m{8}}
{\color{sub}\footnotesize \tF{7} \yr{\textbf{82 Bh}, 82 Ba, 81 Ba, 79 Ba, 75 Ash, \textbf{73 Ch}, 73 Shr}}
\sbs{%
\lead{Role in economic development / why promote}
\begin{itemize}
  \item \textbf{Employment} and self-employment.
  \item \textbf{Capital formation}: mobilizes idle savings.
  \item \textbf{Innovation}, new products, technology.
  \item \textbf{Balanced regional development}; decentralization.
  \item \textbf{Equitable income distribution}.
  \item \textbf{Exports}, import substitution $\rightarrow$ foreign exchange.
  \item Better use of local resources; higher GDP and tax revenue.
  \item Raises living standards; inspires others.
\end{itemize}
\lead{Why Nepal and its government must promote it}
\begin{itemize}
  \item Over a million youths leave for foreign jobs; remittance-dependent economy.
  \item Huge trade deficit; low industrial base.
  \item Unused resources: hydropower, tourism, agriculture, herbs, IT.
  \item Tools: startup loans, incubation, tax holidays, one-window registration, skills
    training, access to credit.
\end{itemize}}{%
\lead{Creativity alone is not enough: a conducive environment (79 Ba, 73 Shr)}
\begin{itemize}
  \item \textbf{Political stability} and consistent policy.
  \item \textbf{Access to finance}: banks, venture capital, collateral-free loans.
  \item \textbf{Infrastructure}: power, roads, internet.
  \item \textbf{Easy registration}, low red tape, fair tax.
  \item \textbf{Market access}, skilled labour, education, incubators.
  \item \textbf{Culture} that accepts failure and risk.
\end{itemize}
\lead{Significance of law enforcement (73 Ch)}
\begin{itemize}
  \item Protects property and \textbf{intellectual property}.
  \item Enforces \textbf{contracts} $\rightarrow$ trust in business deals.
  \item Fair competition; controls cartels and corruption.
  \item Investor confidence $\rightarrow$ capital flows in.
\end{itemize}}

\Q{Entrepreneurship in Nepal's IT sector; would you start a startup; risks and challenges; Silicon Valley}{\m{3}\m{5}\m{8}}
{\color{sub}\footnotesize \tP{3} \yr{76 Ash, \textbf{74 Ch}, \textbf{69 Ch}}}
\sbs{%
\lead{Role in Nepal's IT sector (76 Ash)}
\begin{itemize}
  \item IT services exports (outsourcing, software) earn foreign currency with little
    capital.
  \item Jobs for graduates at home, slowing brain drain.
  \item Digital services: payments (eSewa, Khalti), ride sharing (Pathao), e-commerce
    (Daraz), ed-tech.
  \item Solutions for local problems in agriculture, health, government.
\end{itemize}
\lead{Would I start one after graduation? (74 Ch)}
\begin{itemize}
  \item \textbf{Yes, with preparation}: gain 2--3 years' industry experience and savings,
    validate the idea, find a co-founder, start lean.
  \item Or \textbf{no}: if family needs steady income or the idea is untested. Either answer
    works if justified.
\end{itemize}}{%
\lead{Risks and challenges for an aspiring entrepreneur in Nepal}
\begin{itemize}
  \item Small domestic market; low purchasing power.
  \item Scarce seed and venture capital; banks want collateral.
  \item Restrictions on international payments and foreign investment in small amounts.
  \item Policy instability, bureaucracy, strikes.
  \item Talent loss to migration; social stigma of failure.
  \item Weak IP protection; infrastructure gaps outside cities.
\end{itemize}
\lead{Silicon Valley: the best example (69 Ch)}
\begin{itemize}
  \item Stanford University research + venture capital + risk-taking culture + skilled
    migrants + network effects.
  \item Produced Apple, Google, Intel, HP; failure accepted as learning.
  \item Lesson for Nepal: link universities, investors and startups in clusters.
\end{itemize}}

\penalty0\vspace{2pt}

%=======================================================================
\T{3.5 PYQ Mapping}

\lead{Theory questions, covered in this document}
\begin{itemize}
  \item \textbf{Define motivation; human needs; how needs motivate; intrinsic motivation}
    $\rightarrow$ $\S$3.1. 16 papers, the opener to most theory questions.
  \item \textbf{Techniques of motivation; attitude, group, executive motivation; role of
    management} $\rightarrow$ $\S$3.1. 9 papers.
  \item \textbf{Maslow} (features, satisfaction progression, Devkota 69 Ch, falling from
    grace 78 Bh) $\rightarrow$ $\S$3.2. 9 papers.
  \item \textbf{Maslow vs ERG; ERG and Herzberg vs Maslow; satisfaction-progression vs
    frustration-regression} $\rightarrow$ $\S$3.2. \yr{82 Ba \m{6}, \textbf{78 Bh} \m{3}, \textbf{75 Ch} \m{4},
    71 Shr \m{5}, \textbf{69 Ch} \m{4}, \textbf{68 Ch} \m{5}}
  \item \textbf{McGregor X-Y} \yr{\textbf{81 Bh} \m{5}, 81 Ba \m{5}, 79 Ba \m{6}, 73 Shr \m{5},
    \textbf{72 Ch} \m{8}, \textbf{65 Shr} \m{5}} $\rightarrow$ $\S$3.2.
  \item \textbf{Herzberg} (10 papers) and \textbf{Vroom} (7 papers) $\rightarrow$ $\S$3.2.
  \item \textbf{Fear and punishment} \yr{\textbf{82 Bh} \m{4}}; \textbf{McClelland}: syllabus, never
    set $\rightarrow$ $\S$3.2.
  \item \textbf{Leadership: definition, need, qualities} $\rightarrow$ $\S$3.3. 12 papers.
  \item \textbf{Styles; best style for engineering project / industry / Nepal / start-ups;
    participative management} $\rightarrow$ $\S$3.3. 11 papers.
  \item \textbf{Managerial grid} \yr{\textbf{82 Bh} \m{4}, \textbf{80 Bh} \m{4}, 79 Ba \m{4}, 76 Ash \m{8},
    75 Ash \m{8}, \textbf{70 Ch} \m{5}} $\rightarrow$ $\S$3.3.
  \item \textbf{Leadership approaches: trait, behavioural, contingency, integrated}
    $\rightarrow$ $\S$3.3. 7 papers.
  \item \textbf{Leader vs manager / boss; reader is leader; modern challenges; career
    vision} $\rightarrow$ $\S$3.3. \yr{\textbf{75 Ch}, 75 Ash, \textbf{73 Ch}, \textbf{71 Ch}}
  \item \textbf{Authority and power} \yr{\textbf{66 Bh} \m{5}} $\rightarrow$ $\S$3.3.
  \item \textbf{Entrepreneurship: meaning, characteristics} (8 papers); \textbf{small scale
    unit steps} \yr{\textbf{81 Bh} \m{5}, \textbf{79 Bh} \m{4}, 71 Shr \m{5}} $\rightarrow$ $\S$3.4.
  \item \textbf{Promotion of entrepreneurship, economic role, Nepal, environment, law}
    (7 papers); \textbf{IT sector, startups, Silicon Valley} \yr{76 Ash, \textbf{74 Ch}, \textbf{69 Ch}}
    $\rightarrow$ $\S$3.4.
\end{itemize}

\lead{Short-note prompts from the Detailed PYQ's ``Extra'' section, placed here}
\begin{itemize}
  \item Leadership style \yr{\textbf{68 Ba}}; contingency approach of leadership \yr{\textbf{66 Bh}};
    satisfaction-progression vs frustration-regression \yr{\textbf{75 Ch}, \textbf{69 Ch}}; entrepreneurial
    characteristics \yr{\textbf{79 Bh}}.
  \item 81 Bh's ``State Maslow's need theory with examples [3]'' sits in Ch.\,1 of the
    Detailed PYQ; answered here.
\end{itemize}

\lead{Numerical content}
\begin{itemize}
  \item \textbf{This chapter has no numerical component.} Vroom's $M = E \times I \times V$ is
    the only formula, used qualitatively.
\end{itemize}

\lead{Corrections to the source notes}
\begin{itemize}
  \item AJ Sir slide 10: Maslow is 1943 (``A Theory of Human Motivation''), not ``the
    1940s''; slide 25: the grid's co-author is \textbf{Jane} Mouton, not ``Jake''.
  \item Adhikari p213: fear and punishment slide is criminal deterrence (workplace version
    above); p249: McClelland was at Harvard, not ``Horward''.
\end{itemize}
